\section{The adaptive two-tier system}\label{sec:adaptive}
\subsection{Pass 1: deterministic, from the weights alone}
\begin{enumerate}[label=\arabic*.,leftmargin=1.6em,itemsep=1pt,topsep=1pt]
\item \emph{Forward state (Schr\"odinger).} Run the Gaussian-part sweep for $m_l$ and $S^u_l$, and record each
layer's cut marginals: the level-$1$ marginals $\Phi(\alpha_{l,j})$ and the level-$2$ kernel through the
bivariate orthant law. This gives the channel data $(W_l,K_l)$.
\item \emph{Spectral diagnostics.}
\begin{itemize}[leftmargin=1.2em,itemsep=0pt,topsep=1pt]
\item the gaps $g_l=2\overline{\Phi(\alpha_l)^2}$, computed from the realised $\alpha_{l,j}$ rather than the trajectory;
\item the localisation radii $\ell_k(\eta)$ (Cor.~\ref{cor:loc});
\item the flatness defects $\delta_l=\gamma_l\T\Pi_\perp(F_l-I)\Pi_\perp\gamma_l/\|\gamma_l\|^2$ along the collective gradients $\gamma_l$
(written $g_l$ in stage~18, Thm.~4.2; renamed here to avoid a clash with the gap);
\item the undecided cut $\mathcal C_\theta=\{k:\varphi(\alpha_{L,k})\ge\theta\}$.
\end{itemize}
\item \emph{Backward observables (Heisenberg).} Run the carried-vector recursion for the collective traces,
and form the coherent function $\Ybar(a)$.
\item \emph{Per-neuron diagonals.} Compute $\mathbb E_D[W_L\T S^u_{L-1}W_L]$ and the cubic readout, with history
restricted to the window $[L-\ell_3,L]$ by backward composition.
\item \emph{First output.} Assemble $\hat y^{(1)}=\mathcal J(\cdots)$ (Prop.~\ref{prop:side}).
\end{enumerate}

\subsection{What the first output can reveal: four exact diagnostics}
\begin{proposition}[Wick-orthogonality test]\label{prop:wo}
Let $\eta_k=y_k-\Ybar(a_k)$ be the true incoherent residuals. For every function $\psi$,
$\E[\eta_k\psi(a_k)\mid\mathcal F]=0$. Since the $w_k$ are i.i.d., the statistic $T_\psi=\frac1n\sum_k\eta_k\psi(a_k)$ has mean zero and
variance $\mathrm{Var}(\eta\psi)/n$.
\end{proposition}
\begin{proof}
Stage~18, Theorem~2.1(a): $\E[\eta_k\mid a_k,\mathcal F]=0$. Independence over $k$ gives the variance.
\end{proof}

Apply $T_\psi$ to the first output's residuals $\hat y^{(1)}_k-\Ybar(a_k)$, with $\psi\in\{1,\mathrm{He}_1(a/s),\mathrm{He}_2(a/s)\}$. A value
many standard errors from zero means coherent content has leaked into the per-neuron corrections, which is
double counting between the tiers. The $n$ output neurons are $n$ independent probes of one function $h_Z$,
so this ``uniqueness'' property can be checked from the output itself.

The other three diagnostics:
\begin{itemize}[leftmargin=1.2em,itemsep=1pt,topsep=1pt]
\item[(D2)] \emph{Flatness flags.} Layers with large $\delta_l$ are where the leading-order collective description
is least accurate. They get the exact Wick-order-$2$ couplings $q_l$ and $\rho_l$.
\item[(D3)] \emph{Collective cumulants.} This is stage~19's pass~2. One orthogonal block, read as trace
estimation (Thm.~\ref{thm:hutch}), gives $\hat\kappa_2$ and $\hat\kappa_3$ per layer, with test statistics $D_l$.
\item[(D4)] \emph{Cut-localised orders.} By stage~20, Theorem~2.2, neuron $k$'s order-$r$ correction is at most
$\sigma_k\varphi(\alpha_k)|\mathrm{He}_{r-2}(\alpha_k)||c_{r,k}|$. Only neurons in the undecided cut $\mathcal C_\theta$ need orders beyond $G$.
\end{itemize}

\subsection{Pass 2: dynamic, restarting from a chosen intermediate distribution}
\begin{theorem}[Tier-reduced restart]\label{thm:restart}
Fix a layer $b^*$. Replace the law at $b^*$ by the \emph{tier-reduced state}: the exact Gaussian part
$(m_{b^*},S^u_{b^*})$, together with the collective law (cocycle cumulants and angular spread). In other
words, erase the incoherent cumulants of order $\ge3$ born before $b^*$. Then propagate exactly from $b^*$ to
$L$, with births, cherries and readouts. At leading order, the per-neuron error from the erased history has
ensemble mean square
\[
\mathcal E(b^*)=\Big(\frac{a_3}{6}\Big)^2\sum_{b<b^*}\beta_b^2\prod_{b<l<L}g_l^3,\qquad \beta_b^2=15\,s^6\sum_iB_{b,i}^2 .
\]
Here $B_{b,i}$ are the diagonal and cherry birth amplitudes at layer $b$, computed in pass~1 at cost $O(n)$ per
layer.
\end{theorem}
\begin{proof}
The erased part of neuron $k$'s cubic readout is $\sum_{b<b^*}\sum_i(w_k\T T_{L-1\leftarrow b}e_i)^3B_{b,i}$. By stage~14's
coherence dichotomy (path orthogonality), contributions from different birth layers are uncorrelated over
$W_L$ at leading order. Within a layer, $\E(w\T t)^6=15s^6\|t\|^6$, and
$\|T_{L-1\leftarrow b}e_i\|^6=\prod_{b<l<L}g_l^3\,(1+O(n^{-1/2}))$ by Theorem~\ref{thm:graded} with $k=3$, since $e_i$ is incoherent.
Multiply by the readout coefficient $a_3/6$ of stage~17's Theorem~1.1.
\end{proof}

This is the precise form of ``feeding in a given distribution at an intermediate layer, chosen from the
first attempt''. The layer is $b^*=\max\{b:\mathcal E(b)\le\theta^2\cdot\text{target}^2\}$, read off from pass-1 quantities. The
distribution is the tier-reduced state, which provably carries everything older that still matters.
Everything after $b^*$ is recomputed exactly.

\noindent\textbf{Iteration.}
\begin{enumerate}[label=\roman*.,leftmargin=1.6em,itemsep=0pt,topsep=1pt]
\item Restart at $b^*$, and recompute the window history ($2(L-b^*)$ products).
\item Insert the exact couplings at the flagged layers (D2) and the sampled collective cumulants (D3).
\item Add orders $r\ge3$ only on $\mathcal C_\theta$ (D4).
\item Re-test (Prop.~\ref{prop:wo}, $D_l$). Each refinement adds an orthogonal Wick component, so by stage~18,
Theorem~2.1(c), the ledger is non-increasing. Stop when every diagnostic passes and the ledger is below
target.
\end{enumerate}
The cost is pass~1 (a few hundred products), plus a window of $2\ell_3\approx20$--$28$ products, plus $O(n^3)$ per flagged
layer, plus about $1.5\%$ of the budget for one probe block.

\section*{Status}
\noindent\textbf{Exact:} Theorems~\ref{thm:price} and \ref{thm:hutch}; Propositions~\ref{prop:kraus}, \ref{prop:tp},
\ref{prop:szegedy}, \ref{prop:lift} and \ref{prop:wo}; Theorem~\ref{thm:graded}(a) and (c).

\noindent\textbf{Leading order:} Theorem~\ref{thm:graded}(b), checked by Monte Carlo; Corollary~\ref{cor:loc};
Proposition~\ref{prop:side}; Theorem~\ref{thm:restart}.

\noindent\textbf{Open:}
\begin{itemize}[leftmargin=1.2em,itemsep=0pt,topsep=1pt]
\item the joint-cut correction, meaning the difference between the true multi-layer cut law and the
product of marginal channels, at the next order;
\item spectral independence for the cut law, which would give the register walk its own gap;
\item the general-depth Onsager pairing;
\item validation at $n=1024$ when experiments resume.
\end{itemize}
